\section*{अध्याय \ref{ch:b1:complexation} --- संकुलन}

\begin{solution}{exo:b1:complexation:1}
$[\ce{Cu(NH3)4}]^{2+}$: \ce{Cu^2+}, चार \ce{NH3}, N द्वारा बँधे।
$[\ce{Fe(CN)6}]^{4-}$: \ce{Fe^2+}, छह \ce{CN-}, C द्वारा बँधे;
$[\ce{Fe(CN)6}]^{3-}$: वही, \ce{Fe^3+} के साथ। $[\ce{Ag(NH3)2}]^{+}$:
\ce{Ag+}, दो \ce{NH3}, N द्वारा। $[\ce{FeSCN}]^{2+}$: \ce{Fe^3+}, एक
\ce{SCN-}, जो सूत्र के लिखने के ढंग के बावजूद N द्वारा बँधा है।
$[\ce{Zn(OH)4}]^{2-}$: \ce{Zn^2+}, चार \ce{OH-}, O द्वारा।
$[\ce{CaY}]^{2-}$: \ce{Ca^2+}, एक ईडीटीए आयन, दो N और चार O द्वारा।
\end{solution}

\begin{solution}{exo:b1:complexation:2}
$\log K_i = \log\beta_i - \log\beta_{i-1}$: 4.10, 3.41, 2.82, 2.03, और
$\mathrm{p}K_{d,i} = \log K_i$ वही मान लेता है। अंतिम मान इस अभिक्रिया से संबंधित है:
$[\ce{Cu(NH3)4}]^{2+} \rightleftharpoons [\ce{Cu(NH3)3}]^{2+} + \ce{NH3}$।
\end{solution}

\begin{solution}{exo:b1:complexation:3}
क्षेत्र: pNH3 2.03 से नीचे $[\ce{Cu(NH3)4}]^{2+}$, फिर 2.82 तक $n = 3$,
3.41 तक $n = 2$, 4.10 तक $n = 1$, ऊपर \ce{Cu^2+}। pNH3 = 0:
$[\ce{Cu(NH3)4}]^{2+}$; 2.5: $[\ce{Cu(NH3)3}]^{2+}$; 3.5:
$[\ce{CuNH3}]^{2+}$; 5.0: \ce{Cu^2+}।
\end{solution}

\begin{solution}{exo:b1:complexation:4}
$[\ce{FeSCN^2+}]/[\ce{Fe^3+}] = \beta_1[\ce{SCN-}] = 10^{2.96} \times 0.10 =
91$: संकुलित भिन्न $91/92 = \qty{98.9}{\percent}$ है।
\end{solution}

\begin{solution}{exo:b1:complexation:5}
हर $[\mathrm{ML}_k] = \beta_k[\mathrm{M}][\mathrm{L}]^k$; $k$ पर योग
(कुल धातु) से भाग देने पर $[\mathrm{M}]$ हट जाता है। तीन
स्पीशीज़ के साथ $f_i = 1/(1 + [\mathrm{ML}_{i-1}]/[\mathrm{ML}_i] +
[\mathrm{ML}_{i+1}]/[\mathrm{ML}_i]) = 1/(1 + 1/(K_i x) + K_{i+1}x)$, जहाँ
$x = [\mathrm{L}]$। हर तब सबसे छोटा होता है जब उसका अवकलज
$-1/(K_i x^2) + K_{i+1}$ शून्य हो, $x^2 = 1/(K_iK_{i+1})$, जहाँ दोनों
पड़ोसी बराबर हैं, और $\mathrm{pL} = \frac12(\log K_i + \log K_{i+1})$।
$[\ce{Cu(NH3)2}]^{2+}$ के लिए: $\frac12(3.41 + 2.82) = 3.12$।
\end{solution}

\begin{solution}{exo:b1:complexation:6}
$[\ce{FeSCN}]^{2+} + \ce{F-} \rightleftharpoons [\ce{FeF}]^{2+} + \ce{SCN-}$,
$K = 10^{6.85 - 2.96} = 10^{3.89}$। अनुपात $= K[\ce{F-}]/[\ce{SCN-}] =
10^{3.89} \times 0.010/0.10 = 780$: लाल \omterm{def:b1:complexation:complex}{संकुल} लगभग पूरी तरह
प्रतिस्थापित हो जाता है और रंग फीका पड़ जाता है।
\end{solution}

\begin{solution}{exo:b1:complexation:7}
$[\ce{MgY}]^{2-} + \ce{Ca^2+} \rightleftharpoons [\ce{CaY}]^{2-} + \ce{Mg^2+}$,
$K = 10^{12.69 - 10.90} = 10^{1.79} = 62$। बराबर मात्राओं के साथ
$x^2/(1 - x)^2 = 62$, $x/(1 - x) = 7.9$, $x = 0.89$: मैग्नीशियम का
\qty{89}{\percent} मुक्त हो जाता है।
\end{solution}

\begin{solution}{exo:b1:complexation:8}
\ce{Ag2O(s) + H2O + 4NH3 <=> 2[Ag(NH3)2]+ + 2OH-}, $K = \beta_2^2 \times
10^{-15.43} = 10^{14.44 - 15.43} = 10^{-0.99}$। स्थिरांक
छोटा नहीं है: जैसे ही मुक्त अमोनिया कुछ दशांश मोल प्रति लीटर तक पहुँचती है,
पहली बूँदों से बना भूरा ऑक्साइड (जब अमोनिया क्षारक का काम कर रही थी)
सिल्वर ऐम्मीन में घुल जाता है।
\end{solution}

\begin{solution}{exo:b1:complexation:9}
pH 10: $\alpha = 1 + 10^{1.24} + 10^{-1.96} = 18.4$, $\log\beta' = 12.69 -
1.26 = 11.43$। pH 7: $\alpha = 1 + 10^{4.24} + 10^{4.04} + 10^{0.19} +
\cdots = 10^{4.45}$, $\log\beta' = 8.24$। pH 7 पर सशर्त स्थिरांक
$10^{10}$ से कम है: ईडीटीए के साथ कैल्सियम की अभिक्रिया अनुमापन के लिए
पर्याप्त पूर्ण नहीं है, क्योंकि ईडीटीए का अधिकांश भाग प्रोटॉनित है।
\end{solution}

\begin{solution}{exo:b1:complexation:10}
1.~\ce{CuO(s) + H2O + 4NH3 <=> [Cu(NH3)4]^2+ + 2OH-}, $K = \beta_4 \times
10^{-20.64} = 10^{12.36 - 20.64} = 10^{-8.28}$।
2.~$\mathrm{pH} = 7 + \frac12(9.25 + \log 1.0) = 11.63$, $[\ce{OH-}] =
10^{-2.37}$, और $[\ce{Cu(NH3)4^2+}] = 10^{-8.28}/10^{-4.75} =
\qty{3.0e-4}{mol/L}$: बहुत कम घुलता है।
3.~$\mathrm{pH} = 9.25$, $[\ce{OH-}] = 10^{-4.75}$, और सूत्र
$10^{-8.28 + 9.50} = 10^{1.22}$ देता है, जो अमोनिया की वहन-क्षमता से कहीं अधिक है:
ऑक्साइड तब तक घुलता है जब तक अमोनिया, न कि साम्य, समाप्त न हो जाए।
अमोनियम आयन घुलन से मुक्त हाइड्रॉक्साइड को खपा लेते हैं; इसीलिए
अमोनियम लवण अमोनिया को कॉपर यौगिक घोलने में सहायता देते हैं।
\end{solution}

\begin{solution}{exo:b1:complexation:11}
1.~अभिक्रिया दूसरे विरचन में से पहला घटाने पर मिलती है: $K = K_2/K_1 =
10^{3.90 - 3.32} = 10^{0.58} = 3.8 > 1$, इसलिए $[\ce{AgNH3}]^{+}$ मुख्य
स्पीशीज़ होने पर स्वयं से अभिक्रिया करता है।
2.~\cref{exo:b1:complexation:5} के अनुसार सबसे बड़ा भिन्न
$\mathrm{pNH_3} = \frac12(3.32 + 3.90) = 3.61$ पर है, जहाँ $\beta_1 x = 10^{-0.29}
= 0.51$ और $\beta_2x^2 = 1.0$: $f = 0.51/(1 + 0.51 + 1.0) = 0.20$। सिल्वर का
\qty{20}{\percent} से अधिक कभी $[\ce{AgNH3}]^{+}$ नहीं होता।
\end{solution}

\begin{solution}{exo:b1:complexation:12}
1.~\ce{CaCO3(s) + Y^4- <=> [CaY]^2- + CO3^2-}, $K = \beta K_s = 10^{12.69 -
8.30} = 10^{4.39}$।
2.~$x^2/(0.10 - x) = 10^{4.39}$ से $x = \qty{0.10}{mol/L}$ मिलता है,
$4 \times 10^{-7}$ के भीतर: सारा ईडीटीए प्रयुक्त हो जाता है, प्रति लीटर \qty{10.0}{g} कैल्सियम
कार्बोनेट।
3.~अम्ल में ईडीटीए प्रोटॉनित हो जाता है ($\mathrm{p}K_{a4} = 11.24$) और उसका
सशर्त स्थिरांक गिर जाता है; कार्बोनेट भी प्रोटॉनित होता है, जो सहायक है,
पर कैल्सियम को थामता \omterm{def:b1:complexation:complex}{संकुल} ही है। pH ऊँचा रखने से
ईडीटीए \ce{Y^4-} के रूप में बना रहता है।
\end{solution}

\begin{solution}{pb:b1:complexation:1}
\textbf{1.} \ce{Ag+ + NH3 <=> [AgNH3]+}, $\beta_1 = [\ce{AgNH3+}]/([\ce{Ag+}][\ce{NH3}])$;
\ce{Ag+ + 2NH3 <=> [Ag(NH3)2]+}, $\beta_2 = [\ce{Ag(NH3)2+}]/([\ce{Ag+}][\ce{NH3}]^2)$।
\textbf{2.} $\log K_1 = 3.32$, $\log K_2 = 7.22 - 3.32 = 3.90$।
\textbf{3.} $[\ce{AgNH3}]^{+}$ pNH3 3.90 और 3.32 के बीच प्रभावी होता,
एक ऐसा अंतराल जो उलटा चलता है: उसका कोई क्षेत्र नहीं है।
\textbf{4.} \ce{2[AgNH3]+ <=> Ag+ + [Ag(NH3)2]+}, $K = K_2/K_1 = 10^{0.58} =
3.8$।
\textbf{5.} नीचे $[\ce{Ag(NH3)2}]^{+}$ और ऊपर \ce{Ag+}, सीमा $\frac12\log\beta_2 = 3.61$ पर।
\textbf{6.} $\beta_2[\ce{NH3}]^2 = 10^{7.22 - 2.00} = 10^{5.22} = \num{1.7e5}$।
\textbf{7.} $\mathrm{p}K_d = \log\beta_2 = 7.22$।
\textbf{8.} \ce{AgCl(s) + 2NH3 <=> [Ag(NH3)2]+ + Cl-}, $K = \beta_2K_s =
10^{7.22 - 9.75} = 10^{-2.53}$।
\textbf{9.} $s = \sqrt{K}\,[\ce{NH3}] = 10^{-1.265} \times 1.0 =
\qty{0.054}{mol/L}$।
\textbf{10.} पानी में $s = 10^{-4.875} = \qty{1.3e-5}{mol/L}$: लगभग 4000
गुना अधिक।
\textbf{11.} $0.0543 \times 143.4 = \qty{7.8}{g}$।
\textbf{12.} $[\ce{Ag+}] = K_s/[\ce{Cl-}] = 10^{-9.75}/0.054 =
\qty{3.3e-9}{mol/L} \ll s$।
\textbf{13.} $s = \sqrt{K}(1.0 - 2s)$, $s = 0.0543/(1 + 2 \times 0.0543) =
\qty{0.049}{mol/L}$।
\textbf{14.} \qty{1.0}{mol/L} अमोनिया का pH 11.6 है, जो $9.25 + 1$ से अधिक है:
उसका केवल \qty{0.4}{\percent} \ce{NH4+} है।
\textbf{15.} $K = 10^{7.22 - 12.27} = 10^{-5.05}$।
\textbf{16.} $s = 10^{-2.525} = \qty{3.0e-3}{mol/L}$, \qty{0.56}{g/L}।
\textbf{17.} $K = 10^{-8.85}$, $s = 10^{-4.425} = \qty{3.8e-5}{mol/L}$,
\qty{8.8}{mg/L}।
\textbf{18.} हर अवक्षेप पर तनु अमोनिया डालिए: केवल क्लोराइड
घुलता है। सांद्र अमोनिया ब्रोमाइड को, धीरे और आंशिक रूप से,
घोलती है; आयोडाइड नहीं घुलता।
\textbf{19.} $s = 1.0/187.8 = \qty{5.3e-3}{mol/L}$; $[\ce{NH3}] = s/\sqrt{K}
= \num{5.3e-3}/10^{-2.525} = \qty{1.8}{mol/L}$।
\textbf{20.} $s = \qty{4.3e-3}{mol/L}$ और $[\ce{NH3}] = \num{4.3e-3}/10^{-4.425}
= \qty{113}{mol/L}$, शुद्ध पानी में पानी की सांद्रता (\qty{55}{mol/L}) का
दोगुना: असंभव। अमोनिया सिल्वर आयोडाइड को नहीं घोल सकती।
\textbf{21.} \ce{[Ag(NH3)2]+ + Cl- + 2H3O+ -> AgCl(s) + 2NH4+ + 2H2O},
$K = 10^{2.53} \times (10^{9.25})^2 = 10^{21.03}$: सफ़ेद अवक्षेप
फिर प्रकट होता है, जो सिल्वर क्लोराइड की पुष्टि करता है।
\textbf{22.} $s = 1.0/143.4 = \qty{6.97e-3}{mol/L}$; $[\ce{NH3}] =
s/\sqrt{K} = \num{6.97e-3}/10^{-1.265} = \qty{0.128}{mol/L}$।
\textbf{23.} $2s = \qty{0.0139}{mol/L}$।
\textbf{24.} $[\ce{Ag+}] = 10^{-9.75}/\num{6.97e-3} = \qty{2.6e-8}{mol/L}$,
जो $s$ की तुलना में नगण्य है।
\textbf{25.} $0.128 + 0.014 = $ \textbf{\qty{0.142}{mol/L} अमोनिया}:
लीटर में \qty{0.142}{mol} ठीक \qty{1.0}{g} सिल्वर
क्लोराइड घोलती है।
\end{solution}
